% GENERATED FILE. Edit canonical lessons, then run npm run generate:books.
% canonical-source-hash: fnv1a64:76c221c8dc766703
% canonical-lessons: SA-C371-vikalpah, SA-C371-pratyavartate, SA-C371-vilambate, SA-C371-avadhih, SA-C371-upakramate, SA-R371-review-364-371

\chapter{Other Options}
\label{ch:sa-other-options}
\hlchaptermodality{371}

\begin{chapteropening}
\textbf{By the end of this chapter:} \emph{I can say what to do if the first plan does not work.}

Complete the last lesson of chapter 371: I can say what to do if the first plan does not work.
\end{chapteropening}

\section[\texorpdfstring{\sk{विकल्पः} (vikalpaḥ)}{विकल्पः (vikalpaḥ)}]{\sk{विकल्पः} (vikalpaḥ) — an option; an alternative}

\label{lesson:SA-C371-vikalpah}



\begin{quote}
Before the new one: say the Sanskrit for favourable, then the Sanskrit for adverse.
\end{quote}

\subsection*{You'll want to know: \sk{विकल्पः}}
\textbf{\sk{विकल्पः}} — \emph{vikalpaḥ} — \textquotedblleft{}an option; an alternative\textquotedblright{}.

\begin{quote}
\sk{यदि} \sk{धूमयानं} \sk{न} \sk{अस्ति}, \sk{तर्हि} \sk{लोकयानं} \sk{विकल्पः} \sk{अस्ति।}
\end{quote}

\begin{quote}
\emph{yadi dhūmayānaṃ na asti, tarhi lokayānaṃ vikalpaḥ asti.}
\end{quote}

\begin{quote}
If there is no train, the bus is an option.
\end{quote}

A masculine noun.

\begin{cousinweb}
\textbf{\sk{विकल्पः}} (\emph{vikalpaḥ}) is built on \emph{vi}, \textquotedblleft{}apart\textquotedblright{}, and the root √kḷp, \textquotedblleft{}to be fit, to arrange\textquotedblright{}; MW gives \emph{vi-kalpa} \textquotedblleft{}alternation, alternative, option\textquotedblright{} (MW).
\end{cousinweb}

\begin{grammarlens}[title={What you've built}]
One more word for saying what happens if.
\end{grammarlens}

\subsection*{Your turn}
\begin{itemize}
\raggedright
  \item \emph{Say it:} \emph{vikalpaḥ}
  \item \emph{Say it:} \emph{vikalpaḥ}, once more
  \item \emph{Say it:} \emph{yadi dhūmayānaṃ na asti, tarhi lokayānaṃ vikalpaḥ asti.} — If there is no train, the bus is an option.
  \item \emph{Recall:} say the Sanskrit for favourable, then the Sanskrit for adverse, then say \emph{vikalpaḥ} again
\end{itemize}

\begin{tcolorbox}[breakable,colback=teal!4,colframe=teal!35!black,title={Before you move on}]
What does \sk{विकल्पः} mean? (\textquotedblleft{}An option; an alternative\textquotedblright{}.) Say it once more.
\end{tcolorbox}
\section[\texorpdfstring{\sk{प्रत्यावर्तते} (pratyāvartate)}{प्रत्यावर्तते (pratyāvartate)}]{\sk{प्रत्यावर्तते} (pratyāvartate) — turns back; returns}

\label{lesson:SA-C371-pratyavartate}



\begin{quote}
Before the new one: say the Sanskrit for adverse, then the Sanskrit for an option.
\end{quote}

\subsection*{You'll want to know: \sk{प्रत्यावर्तते}}
\textbf{\sk{प्रत्यावर्तते}} — \emph{pratyāvartate} — \textquotedblleft{}turns back; returns\textquotedblright{}.

\begin{quote}
\sk{यदि} \sk{मार्गे} \sk{जलम्} \sk{अस्ति}, \sk{तर्हि} \sk{यात्री} \sk{प्रत्यावर्तते।}
\end{quote}

\begin{quote}
\emph{yadi mārge jalam asti, tarhi yātrī pratyāvartate.}
\end{quote}

\begin{quote}
If there is water on the road, the traveller turns back.
\end{quote}

Past (he …): \sk{सः} \sk{ग्रामं} \sk{प्रत्यावृत्तवान्}, \emph{saḥ grāmaṃ pratyāvṛttavān}.

\begin{cousinweb}
\textbf{\sk{प्रत्यावर्तते}} (\emph{pratyāvartate}) is built on \emph{prati}, \textquotedblleft{}back\textquotedblright{}, \emph{ā} and the root √vṛt, \textquotedblleft{}to turn\textquotedblright{}; MW gives \emph{praty-ā-√vṛt} \textquotedblleft{}to return, to come back\textquotedblright{} (MW).
\end{cousinweb}

\begin{grammarlens}[title={What you've built}]
One more word for saying what happens if.
\end{grammarlens}

\subsection*{Your turn}
\begin{itemize}
\raggedright
  \item \emph{Say it:} \emph{pratyāvartate}
  \item \emph{Say it:} \emph{pratyāvartate}, once more
  \item \emph{Say it:} \emph{yadi mārge jalam asti, tarhi yātrī pratyāvartate.} — If there is water on the road, the traveller turns back.
  \item \emph{Recall:} say the Sanskrit for adverse, then the Sanskrit for an option, then say \emph{pratyāvartate} again
\end{itemize}

\begin{tcolorbox}[breakable,colback=teal!4,colframe=teal!35!black,title={Before you move on}]
What does \sk{प्रत्यावर्तते} mean? (\textquotedblleft{}Turns back; returns\textquotedblright{}.) Say it once more.
\end{tcolorbox}
\section[\texorpdfstring{\sk{विलम्बते} (vilambate)}{विलम्बते (vilambate)}]{\sk{विलम्बते} (vilambate) — is late; delays}

\label{lesson:SA-C371-vilambate}



\begin{quote}
Before the new one: say the Sanskrit for an option, then the Sanskrit for turns back.
\end{quote}

\subsection*{You'll want to know: \sk{विलम्बते}}
\textbf{\sk{विलम्बते}} — \emph{vilambate} — \textquotedblleft{}is late; delays\textquotedblright{}.

\begin{quote}
\sk{यदि} \sk{भवान्} \sk{विलम्बते}, \sk{तर्हि} \sk{अहं} \sk{गमिष्यामि।}
\end{quote}

\begin{quote}
\emph{yadi bhavān vilambate, tarhi ahaṃ gamiṣyāmi.}
\end{quote}

\begin{quote}
If you are late, I will go.
\end{quote}

Past (he …): \sk{सः} \sk{विलम्बितवान्}, \emph{saḥ vilambitavān}.

\begin{cousinweb}
\textbf{\sk{विलम्बते}} (\emph{vilambate}) is built on \emph{vi} and the root √lamb, \textquotedblleft{}to hang\textquotedblright{}; MW gives \emph{vi-√lamb} \textquotedblleft{}to hang down\textquotedblright{} and \textquotedblleft{}to linger, to delay, to tarry\textquotedblright{} (MW).
\end{cousinweb}

\begin{grammarlens}[title={What you've built}]
One more word for saying what happens if.
\end{grammarlens}

\subsection*{Your turn}
\begin{itemize}
\raggedright
  \item \emph{Say it:} \emph{vilambate}
  \item \emph{Say it:} \emph{vilambate}, once more
  \item \emph{Say it:} \emph{yadi bhavān vilambate, tarhi ahaṃ gamiṣyāmi.} — If you are late, I will go.
  \item \emph{Recall:} say the Sanskrit for an option, then the Sanskrit for turns back, then say \emph{vilambate} again
\end{itemize}

\begin{tcolorbox}[breakable,colback=teal!4,colframe=teal!35!black,title={Before you move on}]
What does \sk{विलम्बते} mean? (\textquotedblleft{}Is late; delays\textquotedblright{}.) Say it once more.
\end{tcolorbox}
\section[\texorpdfstring{\sk{अवधिः} (avadhiḥ)}{अवधिः (avadhiḥ)}]{\sk{अवधिः} (avadhiḥ) — a time limit; a term}

\label{lesson:SA-C371-avadhih}



\begin{quote}
Before the new one: say the Sanskrit for turns back, then the Sanskrit for is late.
\end{quote}

\subsection*{You'll want to know: \sk{अवधिः}}
\textbf{\sk{अवधिः}} — \emph{avadhiḥ} — \textquotedblleft{}a time limit; a term\textquotedblright{}.

\begin{quote}
\sk{यदि} \sk{अवधिः} \sk{श्वः} \sk{अस्ति}, \sk{तर्हि} \sk{अद्य} \sk{एव} \sk{कार्यं} \sk{करोतु।}
\end{quote}

\begin{quote}
\emph{yadi avadhiḥ śvaḥ asti, tarhi adya eva kāryaṃ karotu.}
\end{quote}

\begin{quote}
If the deadline is tomorrow, do the work today.
\end{quote}

A masculine noun.

\begin{cousinweb}
\textbf{\sk{अवधिः}} (\emph{avadhiḥ}) is built on \emph{ava}, \textquotedblleft{}down\textquotedblright{}, and the root √dhā, \textquotedblleft{}to put\textquotedblright{}; MW gives \emph{ava-dhi} \textquotedblleft{}a term, a limit\textquotedblright{} and \textquotedblleft{}a conclusion\textquotedblright{} (MW).
\end{cousinweb}

\begin{grammarlens}[title={What you've built}]
One more word for saying what happens if.
\end{grammarlens}

\subsection*{Your turn}
\begin{itemize}
\raggedright
  \item \emph{Say it:} \emph{avadhiḥ}
  \item \emph{Say it:} \emph{avadhiḥ}, once more
  \item \emph{Say it:} \emph{yadi avadhiḥ śvaḥ asti, tarhi adya eva kāryaṃ karotu.} — If the deadline is tomorrow, do the work today.
  \item \emph{Recall:} say the Sanskrit for turns back, then the Sanskrit for is late, then say \emph{avadhiḥ} again
\end{itemize}

\begin{tcolorbox}[breakable,colback=teal!4,colframe=teal!35!black,title={Before you move on}]
What does \sk{अवधिः} mean? (\textquotedblleft{}A time limit; a term\textquotedblright{}.) Say it once more.
\end{tcolorbox}
\section[\texorpdfstring{\sk{उपक्रमते} (upakramate)}{उपक्रमते (upakramate)}]{\sk{उपक्रमते} (upakramate) — sets about; undertakes}

\label{lesson:SA-C371-upakramate}



\begin{quote}
Before the new one: say the Sanskrit for is late, then the Sanskrit for a time limit.
\end{quote}

\subsection*{You'll want to know: \sk{उपक्रमते}}
\textbf{\sk{उपक्रमते}} — \emph{upakramate} — \textquotedblleft{}sets about; undertakes\textquotedblright{}.

\begin{quote}
\sk{यदि} \sk{योजना} \sk{सज्जा} \sk{अस्ति}, \sk{तर्हि} \sk{सः} \sk{कार्यम्} \sk{उपक्रमते।}
\end{quote}

\begin{quote}
\emph{yadi yojanā sajjā asti, tarhi saḥ kāryam upakramate.}
\end{quote}

\begin{quote}
If the plan is ready, he sets about the work.
\end{quote}

Past (he …): \sk{सः} \sk{कार्यम्} \sk{उपक्रान्तवान्}, \emph{saḥ kāryam upakrāntavān}.

\begin{cousinweb}
\textbf{\sk{उपक्रमते}} (\emph{upakramate}) is built on \emph{upa}, \textquotedblleft{}towards\textquotedblright{}, and the root √kram, \textquotedblleft{}to step\textquotedblright{}; MW gives \emph{upa-√kram} \textquotedblleft{}to approach\textquotedblright{} and \textquotedblleft{}to set about, to undertake, to begin\textquotedblright{} (MW).
\end{cousinweb}

\begin{grammarlens}[title={What you've built}]
One more word for saying what happens if.
\end{grammarlens}

\subsection*{Your turn}
\begin{itemize}
\raggedright
  \item \emph{Say it:} \emph{upakramate}
  \item \emph{Say it:} \emph{upakramate}, once more
  \item \emph{Say it:} \emph{yadi yojanā sajjā asti, tarhi saḥ kāryam upakramate.} — If the plan is ready, he sets about the work.
  \item \emph{Recall:} say the Sanskrit for is late, then the Sanskrit for a time limit, then say \emph{upakramate} again
\end{itemize}

\begin{tcolorbox}[breakable,colback=teal!4,colframe=teal!35!black,title={Before you move on}]
What does \sk{उपक्रमते} mean? (\textquotedblleft{}Sets about; undertakes\textquotedblright{}.) Say it once more.
\end{tcolorbox}
\section[(dialogue)]{Looking back: from Ask an Expert to Other Options}

\label{lesson:SA-R371-review-364-371}



\begin{quote}
Say the Sanskrit for an expert and deliberation.
\end{quote}

\subsection*{Your turn}
Say these aloud:

\begin{itemize}
\raggedright
  \item an expert, deliberation, a commentary, an essay, investigates
  \item a principle, reality, wisdom, curiosity, notices
  \item a conclusion, a review, a calculation, considers, ascertains
  \item a possibility, thinks likely, for certain, swiftly, expects
  \item a misgiving, conjectures, hopes for, an inference, so they say
  \item luck, fate, by accident, comes by, a likelihood
  \item a situation, a state of things, turns into, favourable, adverse
  \item an option, turns back, is late, a time limit, sets about
\end{itemize}

\begin{tcolorbox}[breakable,colback=teal!4,colframe=teal!35!black,title={Before you move on}]
An expert? (\textbf{\sk{विशेषज्ञः}}, \emph{viśeṣajñaḥ}.) Deliberation? (\textbf{\sk{विमर्शः}}, \emph{vimarśaḥ}.) A commentary? (\textbf{\sk{टीका}}, \emph{ṭīkā}.) An essay? (\textbf{\sk{निबन्धः}}, \emph{nibandhaḥ}.) Investigates? (\textbf{\sk{निरूपयति}}, \emph{nirūpayati}.) A principle? (\textbf{\sk{सिद्धान्तः}}, \emph{siddhāntaḥ}.) Reality? (\textbf{\sk{तत्त्वम्}}, \emph{tattvam}.) Wisdom? (\textbf{\sk{प्रज्ञा}}, \emph{prajñā}.) Curiosity? (\textbf{\sk{जिज्ञासा}}, \emph{jijñāsā}.) Notices? (\textbf{\sk{उपलक्षयति}}, \emph{upalakṣayati}.) A conclusion? (\textbf{\sk{निष्कर्षः}}, \emph{niṣkarṣaḥ}.) A review? (\textbf{\sk{समीक्षा}}, \emph{samīkṣā}.) A calculation? (\textbf{\sk{गणना}}, \emph{gaṇanā}.) Considers? (\textbf{\sk{विचारयति}}, \emph{vicārayati}.) Ascertains? (\textbf{\sk{निश्चिनोति}}, \emph{niścinoti}.) A possibility? (\textbf{\sk{सम्भवः}}, \emph{sambhavaḥ}.) Thinks likely? (\textbf{\sk{सम्भावयति}}, \emph{sambhāvayati}.) For certain? (\textbf{\sk{ध्रुवम्}}, \emph{dhruvam}.) Swiftly? (\textbf{\sk{क्षिप्रम्}}, \emph{kṣipram}.) Expects? (\textbf{\sk{अपेक्षते}}, \emph{apekṣate}.) A misgiving? (\textbf{\sk{शंका}}, \emph{śaṅkā}.) Conjectures? (\textbf{\sk{तर्कयति}}, \emph{tarkayati}.) Hopes for? (\textbf{\sk{आशंसते}}, \emph{āśaṃsate}.) An inference? (\textbf{\sk{अनुमानम्}}, \emph{anumānam}.) So they say? (\textbf{\sk{किल}}, \emph{kila}.) Luck? (\textbf{\sk{भाग्यम्}}, \emph{bhāgyam}.) Fate? (\textbf{\sk{दैवम्}}, \emph{daivam}.) By accident? (\textbf{\sk{अकस्मात्}}, \emph{akasmāt}.) Comes by? (\textbf{\sk{उपलभते}}, \emph{upalabhate}.) A likelihood? (\textbf{\sk{सम्भावना}}, \emph{sambhāvanā}.) A situation? (\textbf{\sk{स्थितिः}}, \emph{sthitiḥ}.) A state of things? (\textbf{\sk{अवस्था}}, \emph{avasthā}.) Turns into? (\textbf{\sk{परिणमति}}, \emph{pariṇamati}.) Favourable? (\textbf{\sk{अनुकूलः}}, \emph{anukūlaḥ}.) Adverse? (\textbf{\sk{प्रतिकूलः}}, \emph{pratikūlaḥ}.) An option? (\textbf{\sk{विकल्पः}}, \emph{vikalpaḥ}.) Turns back? (\textbf{\sk{प्रत्यावर्तते}}, \emph{pratyāvartate}.) Is late? (\textbf{\sk{विलम्बते}}, \emph{vilambate}.) A time limit? (\textbf{\sk{अवधिः}}, \emph{avadhiḥ}.) Sets about? (\textbf{\sk{उपक्रमते}}, \emph{upakramate}.)
\end{tcolorbox}
